\section{Provenance and references}

The target classical theorem was the finite-dimensional
\(C^1\) implicit function theorem described in the user's
linked overview \cite{wikipedia}. The proof here was
developed directly from quantitative rational certificates.
The constructive literature \cite{diener2011,diener2009}
provides context for the role of quantitative uniqueness.
The classical article \cite{oliveira} is a useful
comparison, not a foundational premise of this proof.

The ProveIt sources were inspected at the exact commit
specified in Section~\ref{sec:lean}. Their relevant
interfaces and limitations are documented in
\code{references/proveit/repo\_audit.md}; the accompanying
manifest gives a commit-pinned URL and blob hash for
each copied reference source. These files are evidence
for the reuse assessment and are not part of the
self-contained \code{FiniteCore} build target.

\begingroup
\raggedright
\begin{thebibliography}{9}

\bibitem{diener2011}
H. Diener and P. Schuster,
\emph{Uniqueness, continuity and the existence of implicit
functions in constructive analysis},
LMS Journal of Computation and Mathematics \textbf{14}
(2011), 127--136.
\url{https://doi.org/10.1112/S1461157010000057}.

\bibitem{diener2009}
H. Diener and P. Schuster,
\emph{Uniqueness, Continuity, and Existence of Implicit
Functions in Constructive Analysis},
in CCA 2009, OASIcs \textbf{11}, 131--140.
\url{https://doi.org/10.4230/OASIcs.CCA.2009.2265}.

\bibitem{oliveira}
O. Rio Branco de Oliveira,
\emph{The Implicit and the Inverse Function theorems:
easy proofs},
Real Analysis Exchange \textbf{39}(1) (2013/2014),
207--218; arXiv:1212.2066.
\url{https://arxiv.org/abs/1212.2066}.

\bibitem{proveit}
V. Reshetnikov and contributors,
\emph{ProveIt: Logic/PrimitiveRecursiveArithmetic},
snapshot at commit
\code{d36b35eed5bd01af4ade66bf539622495813ccc8}.
\url{https://github.com/VladimirReshetnikov/ProveIt/tree/d36b35eed5bd01af4ade66bf539622495813ccc8/Logic/PrimitiveRecursiveArithmetic}.
Selected object-theory derivation sources are listed
in the package manifest.

\bibitem{mathlib}
The Mathlib community,
\emph{Mathlib.Analysis.Calculus.ImplicitContDiff},
official documentation, inspected November 21, 2026.
\url{https://leanprover-community.github.io/mathlib4_docs/Mathlib/Analysis/Calculus/ImplicitContDiff.html}.
The documentation identifies
\code{ContDiffAt.implicitFunction} as noncomputable
and gives its smoothness and local inverse properties.

\bibitem{wikipedia}
\emph{Implicit function theorem},
Wikipedia, user-specified overview, inspected
November 21, 2026.
\url{https://en.wikipedia.org/wiki/Implicit_function_theorem}.

\end{thebibliography}
\endgroup
